\subsection{适配序列与相伴微分形式}

在本小节中，直到第 5 小节，我们用 A 表示一个复交换幺 Banach 代数，用 n 表示一个整数 \(\geqslant\) 1。

当我们谈到 \(C^{n}\) 的开子集上的无穷次可微函数时，指的是对于底实流形结构无穷次可微的函数。所用的微分演算概念都是相对于这个结构的。

定义 1. --- 设 \(\boldsymbol{a} = (a_{1},\ldots ,a_{n})\in \mathrm{A}^{n}\)，设 \(h\colon \mathbf{C}^n\to \mathbf{C}\) 是一个映射，并设 \(u_{1},\ldots ,u_{n}\) 是从 \(\mathbf{C}^n\) 到 A 中的映射。我们说序列 \((h,u_1,\dots ,u_n)\) 适配于 \(\boldsymbol{a}\)，如果

\begin{enumerate}
\def\labelenumi{(\roman{enumi})}
\item
  映射 h 是无穷次可微的，具有紧支集，且在 \(\mathrm{Sp}^n (\boldsymbol {a})\) 的一个邻域中等于 1；
\item
  诸映射 \(u_1, \ldots, u_n\) 都是无穷次可微的；
\item
  对每个 \(\boldsymbol{z} = (z_{1},\dots ,z_{n})\in \mathbf{C}^{n}\)，我们有
\end{enumerate}

\[
h (\boldsymbol {z}) + (z _ {1} - a _ {1}) u _ {1} (\boldsymbol {z}) + \dots + (z _ {n} - a _ {n}) u _ {n} (\boldsymbol {z}) = 1.\tag{1}
\]

\(C^{n}\) 上的、以 A 为系数的 2n 次微分形式，由

\[
\omega = \bigwedge_ {i = 1} ^ {n} (d u _ {i} \wedge d z _ {i})
\]

所定义，称为与 \((h, u_1, \ldots, u_n)\) 相伴的微分形式。

若 \((h, u_{1}, \ldots, u_{n})\) 适配于 a，则对 (1) 微分，我们得到等式

\[
d h = - \sum_ {i = 1} ^ {n} u _ {i} d z _ {i} - \sum_ {i = 1} ^ {n} (z _ {i} - a _ {i}) d u _ {i}\tag{2}
\]

由此，对每个满足 \(1 \leqslant i \leqslant n\) 的 i，得到关系式

\[
d h \wedge d z _ {i} \wedge \bigwedge_ {j \neq i} (d u _ {j} \wedge d z _ {j}) = - (z _ {i} - a _ {i}) \omega .\tag{3}
\]

引理 1. --- 设 U 是 \(\mathbf{C}^{n}\) 的一个开子集。设 K 是 U 的一个紧子集。则存在一个从 \(\mathbf{C}^{n}\) 到 C 中的无穷次可微函数 h，它在 K 上等于 1 且具有含于 U 中的紧支集。

设 V 是 K 的一个相对紧开邻域，使得 \(\overline{V}\) 含于 U (TG, I, p.~65, prop. 10)。存在 \(C^{n}\) 到 C 中的一个无穷次可微函数 h，其支集含于 V 且在 K 上等于 1 (VAR, R1, p.~40, 5.3.6)。这个函数具有所要求的性质。

例. --- 假定 n = 1。设 \(a \in A\)。对 \(\mathrm{Sp}(a)\) 的每个开邻域 U，存在一个从 C 到 C 中的无穷次可微函数 h，其紧支集含于 U，且在 \(\mathrm{Sp}(a)\) 的一个邻域中等于 1 (VAR, R1, p.~40, 5.3.6)。令

\[
u (z) = (1 - h (z)) (z - a) ^ {- 1}
\]

（对 \(z \in \mathbf{C} - \mathrm{Sp}(a)\)），并令 \(u(z) = 0\)（若 \(z \in \mathrm{Sp}(a)\)）。数对 \((h, u)\) 适配于 a，而相伴微分形式是 \(\omega = du \wedge dz\)。

引理 2. --- 设 \(\boldsymbol{a} = (a_{1},\ldots ,a_{n})\in \mathrm{A}^{n}\)。存在无穷次可微映射 \(v_{1},\ldots ,v_{n}\)（从 \(\mathbf{C}^n -\mathrm{Sp}^n (\boldsymbol {a})\) 到 A 中），使得

\[
(z _ {1} - a _ {1}) v _ {1} (\boldsymbol {z}) + \dots + (z _ {n} - a _ {n}) v _ {n} (\boldsymbol {z}) = 1
\]

对一切 \(z = (z_{1}, \ldots, z_{n}) \in \mathbf{C}^{n} - \mathrm{Sp}^{n}(\boldsymbol{a})\)。

设 \(\boldsymbol{w} = (w_{1}, \ldots, w_{n}) \in \mathbf{C}^{n} - \operatorname{Sp}^{n}(\boldsymbol{a})\)。依联合谱的定义，存在 A 中的元素 \(b_{1}, \ldots, b_{n}\)，使得

\[
(w _ {1} - a _ {1}) b _ {1} + \dots + (w _ {n} - a _ {n}) b _ {n} = 1
\]

(I, p.~41, 第 6 小节)。存在 w 的一个开邻域 \(W_{w}\)，使得 A 的元素 \((z_{1}-a_{1})b_{1}+\cdots+(z_{n}-a_{n})b_{n}\) 当 \(z=(z_{1},\ldots,z_{n})\) 属于 \(W_{w}\) 时是可逆的。于是，对每个满足 \(1\leqslant j\leqslant n\) 的整数 j 及 \(W_{w}\) 中的每个 z，令

\[
u _ {j} (\boldsymbol {z}) = b _ {j} \left(\sum_ {i = 1} ^ {n} (z _ {i} - a _ {i}) b _ {i}\right) ^ {- 1}.
\]

这样定义的函数 \(u_{1}, u_{2}, \ldots, u_{n}\)（\(W_{w}\) 到 A 中的函数）在 \(W_{w}\) 中是无穷次可微的，且我们有

\[
(z _ {1} - a _ {1}) u _ {1} (\boldsymbol {z}) + \dots + (z _ {n} - a _ {n}) u _ {n} (\boldsymbol {z}) = 1
\]

对 \(W_{w}\) 中的一切 z。

由于族 \((\mathrm{W}_{\boldsymbol{w}})_{\boldsymbol{w}\in\mathbf{C}^{n}-\mathrm{Sp}^{n}(\boldsymbol{a})}\) 是 \(\mathbf{C}^{n}-\mathrm{Sp}^{n}(\boldsymbol{a})\) 的一个开覆盖，存在一个局部有限的开加细 \(\mathscr{W}=(\mathrm{W}_{\lambda})_{\lambda\in\mathrm{L}}\) (TG, I, p.~70, th. 5)，且对每个 \(\lambda\in L\)，存在函数 \(u_{1\lambda},\ldots,u_{n\lambda}\)（取值于 A 中，定义且无穷次可微于 \(W_{\lambda}\) 中），使得 \((z_{1}-a_{1})u_{1\lambda}(\boldsymbol{z})+\cdots+(z_{n}-a_{n})u_{n\lambda}(\boldsymbol{z})=1\) 对一切 \(\boldsymbol{z}\)（位于 \(W_{\lambda}\) 中）成立。设 \((f_{\lambda})_{\lambda \in L}\) 是一个从属于覆盖 \(\mathscr{W}\) 的、由无穷次可微函数组成的单位分解 (VAR, R1, p.~40, 5.3.6)。设 \(i\) 是一个满足 \(1 \leqslant i \leqslant n\) 的整数。对每个 \(\lambda \in L\)，设 \(u_{i\lambda}'\) 是从 \(\mathbf{C}^n - \mathrm{Sp}^n(\boldsymbol{a})\) 到 A 中的映射，它由把函数 \(f_{\lambda}u_{i\lambda}\) 延拓（在 \((\mathbf{C}^n - \mathrm{Sp}^n(\boldsymbol{a})) - W_{\lambda}\) 中取 0）而得到。诸函数 \(u_{i\lambda}'\) 是无穷次可微的。由于族 \((\mathrm{Supp}(u_{i\lambda}')_{\lambda \in L}\) 是局部有限的，函数 \(v_i = \sum_{\lambda \in L} u_{i\lambda}'\) 在 \(\mathbf{C}^n - \mathrm{Sp}^n(\boldsymbol{a})\) 中有定义且无穷次可微。

设 \(z \in \mathbf{C}^n - \mathrm{Sp}^n(\boldsymbol{a})\)。用 L' 表示 \(\lambda \in \mathrm{L}\)（满足 \(z \in \mathrm{W}_{\lambda}\)）所成的有限集。则

\[
\begin{array}{c} \sum_ {i = 1} ^ {n} (z _ {i} - a _ {i}) v _ {i} (\boldsymbol {z}) = \sum_ {\lambda \in \mathrm{L} ^ {\prime}} \sum_ {i = 1} ^ {n} (z _ {i} - a _ {i}) u _ {i \lambda} ^ {\prime} (\boldsymbol {z}) \\ = \sum_ {\lambda \in \mathrm{L} ^ {\prime}} f _ {\lambda} (\boldsymbol {z}) \sum_ {i = 1} ^ {n} (z _ {i} - a _ {i}) u _ {i \lambda} (\boldsymbol {z}) = \bigg (\sum_ {\lambda \in \mathrm{L} ^ {\prime}} f _ {\lambda} (\boldsymbol {z}) \bigg) \cdot 1 = 1. \end{array}
\]

引理 3. --- 设 \(\boldsymbol{a} = (a_{1},\ldots ,a_{n})\in \mathrm{A}^{n}\)。设 \(h\) 是一个从 \(\mathbf{C}^n\) 到 \(\mathbf{C}\) 中的无穷次可微映射，在 \(\mathrm{Sp}^n (\boldsymbol {a})\) 的一个邻域中等于 1 且具有紧支集。存在无穷次可微映射 \(u_{1},\ldots ,u_{n}\)（从 \(\mathbf{C}^n\) 到 A 中的），使得序列 \((h,u_1,\dots,u_n)\) 适配于 \(\boldsymbol{a}\)。

设 \(v_{1},\ldots,v_{n}\) 是从 \(\mathbf{C}^{n}-\mathrm{Sp}^{n}(\boldsymbol{a})\) 到 A 中的无穷次可微映射，使得

\[
\sum_ {j = 1} ^ {n} (z _ {j} - a _ {j}) v _ {j} (\boldsymbol {z}) = 1
\]

（对 \(z\) 属于 \(\mathbf{C}^n - \mathrm{Sp}^n(\boldsymbol{a})\)）（引理 2）。设 \(i\) 是一个满足 \(1 \leqslant i \leqslant n\) 的整数。令 \(u_i(z) = (1 - h(z))v_i(z)\)（若 \(z \in \mathbf{C}^n - \mathrm{Sp}^n(\boldsymbol{a})\)），并令 \(u_i(z) = 0\)（若 \(z \in \mathrm{Sp}^n(\boldsymbol{a})\)）。诸映射 \(u_i\) 在 \(\mathbf{C}^n - \mathrm{Sp}^n(\boldsymbol{a})\) 中无穷次可微，且在 \(\mathrm{Sp}^n(\boldsymbol{a})\) 的一个邻域中为零，因而在 \(\mathbf{C}^n\) 中无穷次可微。等式 (1) 在 \(\mathrm{Sp}^n(\boldsymbol{a})\) 上成立，因为诸函数 \(u_i\) 在 \(\mathrm{Sp}^n(\boldsymbol{a})\) 上为零，且 \(h\) 在 \(\mathrm{Sp}^n(\boldsymbol{a})\) 的一个邻域中等于 1。由构造，它在 \(\mathbf{C}^n - \mathrm{Sp}^n(\boldsymbol{a})\) 上也成立。

引理 4. — 设 \(a \in A^{n}\)。设 \((h, u_{1}, \ldots, u_{n})\) 是一个适配于 a 的序列，\(\omega\) 是相伴微分形式。

\begin{enumerate}
\def\labelenumi{\alph{enumi})}
\tightlist
\item
  对 \(i = 1,2,\ldots,n\)，存在一个微分形式 \(\beta_{i}\)（\(\mathbf{C}^{n}\) 上的、\(n - 1\) 次的、以 A 为系数的微分形式），使得
\end{enumerate}

\[
(z _ {i} - a _ {i}) \omega = d (h \beta_ {i} \wedge d z _ {1} \wedge \dots \wedge d z _ {n});
\]

\begin{enumerate}
\def\labelenumi{\alph{enumi})}
\setcounter{enumi}{1}
\item
  微分形式 \(\omega\) 具有含于 \(h\) 的支集中的紧支集；
\item
  存在一个微分形式 \(\beta\)（\(\mathbf{C}^{n}\) 上的、\(n - 1\) 次的、以 A 为系数的微分形式），使得
\end{enumerate}

\[
(n + 1) h \omega - \omega = d (h \beta \wedge d z _ {1} \wedge \dots \wedge d z _ {n}).
\]

设 \(i\) 是一个满足 \(1 \leqslant i \leqslant n\) 的整数。存在 \(\varepsilon_i \in \{-1, 1\}\)，使得

\[
\varepsilon_ {i} \bigwedge_ {j \neq i} d u _ {j} \wedge d z _ {1} \wedge \dots \wedge d z _ {n} = d z _ {i} \wedge \bigwedge_ {j \neq i} (d u _ {j} \wedge d z _ {j}).
\]

令 \(\beta_{i} = \varepsilon_{i} \bigwedge_{j \neq i} du_{j}\)，于是该公式的左端项是 \(\beta_{i} \wedge dz_{1} \wedge \cdots \wedge dz_{n}\)，且 \(d\beta_{i} = 0\)。这样

\[
d \left(h \beta_ {i} \wedge d z _ {1} \wedge \dots \wedge d z _ {n}\right) = d h \wedge \beta_ {i} \wedge d z _ {1} \wedge \dots \wedge d z _ {n} =
\]

\[
d h \wedge d z _ {i} \wedge \bigwedge_ {j \neq i} (d u _ {j} \wedge d z _ {j}) = (z _ {i} - a _ {i}) \omega ,
\]

这是依公式 (3)，由此得断言 \(a\) )。

由断言 \(a\) ) 及公式 (1)，我们推出关系式

\[
\begin{array}{l} \omega = h \omega + (1 - h) \omega = h \omega + \sum_ {i = 1} ^ {n} (z _ {i} - a _ {i}) u _ {i} \omega \\ \qquad \qquad \qquad = h \omega + \sum_ {i = 1} ^ {n} u _ {i}   d (h \beta_ {i} \wedge d z _ {1} \wedge \dots \wedge d z _ {n}), \end{array}
\]

由此 \(\operatorname{Supp}(\omega) \subset \operatorname{Supp}(h)\)，这就证明了 \(b\) )。

最后，令

\[
\beta = \sum_ {i = 1} ^ {n} \varepsilon_ {i} u _ {i} \bigwedge_ {j \neq i} d u _ {j} = \sum_ {i = 1} ^ {n} u _ {i} \beta_ {i}, \text {   and   } \tau = h \beta d z _ {1} \wedge \dots \wedge d z _ {n}.
\]

我们有 \(d\beta = \sum_{i}du_{i}\wedge \beta_{i}\)，因而

\[
d \beta \wedge d z _ {1} \wedge \dots \wedge d z _ {n} = \sum_ {i} d u _ {i} \wedge d z _ {i} \wedge \bigwedge_ {j \neq i} (d u _ {j} \wedge d z _ {j}) = n \omega .
\]

于是

\[
\begin{array}{l} d \tau = d h \wedge \beta \wedge d z _ {1} \wedge \dots \wedge d z _ {n} + h d \beta \wedge d z _ {1} \wedge \dots \wedge d z _ {n} \\ = \sum_ {i = 1} ^ {n} u _ {i} d h \wedge d z _ {i} \wedge \bigwedge_ {j \neq i} (d u _ {j} \wedge d z _ {j}) + n h \omega \\ = - \sum_ {i = 1} ^ {n} u _ {i} (z _ {i} - a _ {i}) \omega + n h \omega = (h - 1) \omega + n h \omega = (n + 1) h \omega - \omega , \end{array}
\]

给定公式 (3) 与 (1)，由此得 c)。

我们现在来研究微分形式 \(\omega\)（与一个适配于 \(\boldsymbol{a}\) 的序列相伴的）如何作为该序列的函数而变化。我们说两个序列 \((h, u_1, \ldots, u_n)\) 与 \((h', u_1', \ldots, u_n')\)（均适配于 \(\boldsymbol{a}\)）是相联的，如果存在一个微分形式 \(\psi\)（\(n-1\) 次的、\(\mathbf{C}^n\) 上的、以 A 为系数且支包含于 \(h\) 与 \(h'\) 的支集之并中），使得相伴的微分形式 \(\omega\) 与 \(\omega'\) 满足

\[
\omega - \omega^ {\prime} = d (\psi \wedge d z _ {1} \wedge d z _ {2} \wedge \dots \wedge d z _ {n}).
\]

让我们从一个初等变形开始：

引理 5. --- 设 \(\boldsymbol{a} \in \mathrm{A}^{n}\)，设 \((h, u_{1}, \ldots, u_{n})\) 是一个适配于 \(\boldsymbol{a}\) 的序列，并设 \(\omega\) 是相伴微分形式。

设 \(w\) 是一个从 \(\mathbf{C}^n\) 到 \(A\) 中的无穷次可微映射，并设 \(i\) 与 \(j\) 是 1 与 \(n\) 之间的两个不同整数。用下式定义 \(u_1', \ldots, u_n'\)：

\[
\begin{array}{c} u _ {i} ^ {\prime} = u _ {i} + (z _ {j} - a _ {j}) w, \quad u _ {j} ^ {\prime} = u _ {j} - (z _ {i} - a _ {i}) w, \\ u _ {k} ^ {\prime} = u _ {k} \text {for} k \neq i, j. \end{array}
\]

则序列 \((h,u_{1}^{\prime},\ldots,u_{n}^{\prime})\) 适配于 a，且与序列 \((h,u_{1},\ldots,u_{n})\) 相联。

我们记 \(dz = dz_{1} \wedge \cdots \wedge dz_{n}\)。由于

\[
\begin{array}{c} \sum_ {k = 1} ^ {n} (z _ {k} - a _ {k}) u _ {k} ^ {\prime} (\boldsymbol {z}) = \sum_ {k = 1} ^ {n} (z _ {k} - a _ {k}) u _ {k} (\boldsymbol {z}) + w (\boldsymbol {z}) (z _ {j} - a _ {j}) (z _ {i} - a _ {i}) \\ - w (\boldsymbol {z}) (z _ {i} - a _ {i}) (z _ {j} - a _ {j}) = 1 - h (\boldsymbol {z}) \end{array}
\]

对一切 \(z \in C^{n}\)，序列 \((h, u_{1}', \ldots, u_{n}')\) 适配于 a。此外，我们有

\[
\begin{array}{l} d u _ {i} ^ {\prime} \wedge d u _ {j} ^ {\prime} \wedge d z _ {1} \wedge \dots \wedge d z _ {n} = \\ \left(d u _ {i} + w d z _ {j} + (z _ {j} - a _ {j}) d w\right) \wedge \left(d u _ {j} - w d z _ {i} - (z _ {i} - a _ {i}) d w\right) \wedge d \boldsymbol {z} \\ = \left(d u _ {i} \wedge d u _ {j} - (z _ {i} - a _ {i}) d u _ {i} \wedge d w - (z _ {j} - a _ {j}) d u _ {j} \wedge d w\right) \wedge d \boldsymbol {z} \end{array}
\]

因此存在 \(\varepsilon\in\{-1,1\}\)，使得 \(\varepsilon(\omega-\omega')\) 等于

\[
\begin{array}{l} d u _ {i} ^ {\prime} \wedge d u _ {j} ^ {\prime} \wedge \bigwedge_ {k \neq i, j} d u _ {k} ^ {\prime} \wedge d \boldsymbol {z} - d u _ {i} \wedge d u _ {j} \wedge \bigwedge_ {k \neq i, j} d u _ {k} \wedge d \boldsymbol {z} \\ = - \Big ((z _ {i} - a _ {i}) d u _ {i} \wedge d w + (z _ {j} - a _ {j}) d u _ {j} \wedge d w \Big) \wedge \bigwedge_ {k \neq i, j} d u _ {k} \wedge d \boldsymbol {z} \\ = - \Big (\sum_ {k = 1} ^ {n} (z _ {k} - a _ {k})   d u _ {k} \Big) \wedge d w \wedge \bigwedge_ {k \neq i, j} d u _ {k} \wedge d \boldsymbol {z} \end{array}
\]

而把 (2) 考虑在内，它就等于

\[
d h \wedge d w \wedge \bigwedge_ {k \neq i, j} d u _ {k} \wedge d \boldsymbol {z} = d \left(h d w \wedge \bigwedge_ {k \neq i, j} d u _ {k} \wedge d \boldsymbol {z}\right),
\]

由此即得结论。

引理 6. — 设 \(\mathbf{a} \in \mathrm{A}^n\)。所有适配于 \(\mathbf{a}\) 的序列都是相联的。

设 \((h,u_{1},\ldots,u_{n})\) 与 \((h',u_{1}',\ldots,u_{n}')\) 是两个适配于 a 的序列，并用 \(\omega\) 与 \(\omega'\) 表示相伴的微分形式。

我们来定义无穷次可微映射

\[
\begin{array}{c} {w _ {i j} = u _ {i} ^ {\prime} u _ {j} - u _ {i} u _ {j} ^ {\prime},} \\ {s _ {i} = u _ {i} ^ {\prime} h - u _ {i} h ^ {\prime},} \end{array}
\]

\[
\begin{array}{l} 1 \leqslant i \leqslant n, 1 \leqslant j \leqslant n, \\ 1 \leqslant i \leqslant n, \end{array}
\]

于是 \(w_{ji} = -w_{ij}\)，且 \(\operatorname{Supp}(s_i) \subset \operatorname{Supp}(h) \cup \operatorname{Supp}(h')\)。

我们置 \(u_i'' = u_i' - s_i\)，\(\boldsymbol{u} = (u_1, \ldots, u_n)\) 与 \(\boldsymbol{u}'' = (u_1'', \ldots, u_n'')\)。再用 \(\boldsymbol{v}_{ij}\) 表示从 \(\mathbf{C}^n\) 到 \(A^n\) 中的映射，其第 \(i\) 个分量是 \((z_j - a_j) w_{ij}\)，第 \(j\) 个分量是 \((z_i - a_i) w_{ji} = -(z_i - a_i) w_{ij}\)，其余分量均为零。则我们有

\[
\boldsymbol {u} ^ {\prime \prime} = \boldsymbol {u} + \sum_ {i <   j} \boldsymbol {v} _ {i j}.
\]

事实上，对每个满足 \(1 \leqslant k \leqslant n\) 的整数 k，右端的第 k 个分量是

\[
\begin{array}{c} u _ {k} + \sum_ {i = 1} ^ {k - 1} (z _ {i} - a _ {i}) w _ {k i} + \sum_ {j = k + 1} ^ {n} (z _ {j} - a _ {j}) w _ {k j} = u _ {k} + \sum_ {i = 1} ^ {n} (z _ {i} - a _ {i}) w _ {k i} \\ = u _ {k} + u _ {k} ^ {\prime} \sum_ {i = 1} ^ {n} (z _ {i} - a _ {i}) u _ {i} - u _ {k} \sum_ {i = 1} ^ {n} (z _ {i} - a _ {i}) u _ {i} ^ {\prime} \\ = u _ {k} + (1 - h) u _ {k} ^ {\prime} - (1 - h ^ {\prime}) u _ {k} = u _ {k} ^ {\prime} - s _ {k}. \end{array}
\]

由归纳法，我们从引理 5（应用于整数 \(i\) 与 \(j\) 以及映射 \(w_{ij}\)）推出，序列 \((h,u_1'',\ldots,u_n'')\) 适配于 \(\boldsymbol{a}\) 且与 \((h,u_1,\ldots,u_n)\) 相联。设 \(\omega''\) 是与 \((h,u_1'',\ldots,u_n'')\) 相伴的微分形式。由于 \(u_i'' = u_i' - s_i\)，我们有

\[
\begin{array}{r} \omega^ {\prime \prime} - \omega^ {\prime} = d (u _ {1} ^ {\prime} - s _ {1}) \wedge d z _ {1} \wedge \dots \wedge d (u _ {n} ^ {\prime} - s _ {n}) \wedge d z _ {n} - \\ d u _ {1} ^ {\prime} \wedge d z _ {1} \wedge \dots \wedge d u _ {n} ^ {\prime} \wedge d z _ {n}, \end{array}
\]

它可表为如下形式的微分形式的一个线性组合（系数为 1 或 -1）

\[
\xi_ {\mathrm{I} _ {1}, \mathrm{I} _ {2}} = \bigwedge_ {i \in \mathrm{I} _ {1}} d s _ {i} \wedge \bigwedge_ {i \in \mathrm{I} _ {2}} d u _ {i} ^ {\prime} \wedge d z _ {1} \wedge \dots \wedge d z _ {n},
\]

其中 \(I_{1}\)（相应地，\(I_{2}\)）是 \(\{1,\ldots,n\}\) 的一个非空子集（相应地，一个子集）。每个微分形式 \(\xi_{I_{1},I_{2}}\) 也可写成如下形式

\[
d (\widetilde {\psi} \wedge d z _ {1} \wedge \dots \wedge d z _ {n}),
\]

其中微分形式 \(\widetilde{\psi}\) 的支包含于 \(s_i\) 的支集，对一切 \(i \in \mathrm{I}_1\)。由于 \(\mathrm{I}_1\) 非空，这个支集包含于 \(\operatorname{Supp}(h) \cup \operatorname{Supp}(h')\)。因而 \((h, u_1'', \ldots, u_n'')\) 与 \((h', u_1', \ldots, u_n')\) 相联，并且写出如下式子即得引理

\[
\omega - \omega^ {\prime} = (\omega - \omega^ {\prime \prime}) + (\omega^ {\prime \prime} - \omega^ {\prime}).
\]

